\section{Mortalidade}

\subsection{Procedimentos}

Nesta questão foram utilizados os microdados do SIM/DATASUS para óbitos de residentes no Amapá \cite{datasus2026} e as populações projetadas pelo IBGE -- Revisão 2024 \cite{ibge2024}. Os anos analisados foram 2010, 2019, 2021, 2022 e 2024.

Foram calculadas a Taxa Bruta de Mortalidade (TBM), as taxas específicas de mortalidade por sexo e idade (${}_nM_x$), indicadores de mortalidade infantil e perinatal, a estrutura de causas de morte e tábuas de vida abreviadas para 2010 e 2024.

A Taxa Bruta de Mortalidade foi calculada por:

\[
TBM = \frac{D}{P} \times 1000
\]

em que $D$ é o total de óbitos no ano e $P$ é a população projetada no mesmo ano.

As taxas específicas de mortalidade foram calculadas por:

\[
{}_nM_x = \frac{{}_nD_x}{{}_nP_x} \times 1000
\]

em que ${}_nD_x$ representa os óbitos no grupo etário e ${}_nP_x$ representa a população do mesmo grupo. Neste caso, ${}_nM_x$ é apresentado por mil habitantes.

\subsection{Taxa Bruta e Taxas Específicas de Mortalidade}

A Tabela \ref{tab:tbm_q3} apresenta a Taxa Bruta de Mortalidade para os anos selecionados. A TBM aumenta entre 2010 e 2021, com pico em 2021, ano marcado pela pandemia de COVID-19. Em 2022 e 2024, a taxa diminui em relação a 2021, mas permanece acima do valor observado em 2010.

\begin{table}[H]
\centering
\small
\caption{Taxa Bruta de Mortalidade. Amapá, anos selecionados.}
\label{tab:tbm_q3}
\begin{tabular}{rrrr}
\toprule
\textbf{Ano} & \textbf{Óbitos} & \textbf{População projetada} & \textbf{TBM por mil} \\
\midrule
2010 & 2.276 & 686.987 & 3,31 \\
2019 & 3.407 & 779.468 & 4,37 \\
2021 & 4.578 & 790.871 & 5,79 \\
2022 & 3.735 & 795.457 & 4,70 \\
2024 & 3.798 & 802.837 & 4,73 \\
\bottomrule
\end{tabular}
\end{table}

As taxas específicas por sexo e idade mostram o padrão esperado: mortalidade mais elevada no primeiro ano de vida, queda nas idades infantis e aumento progressivo nas idades adultas e idosas. Também se observa maior mortalidade masculina em quase todos os grupos etários, principalmente nas idades jovens e adultas.

\begin{figure}[H]
    \centering
    \caption{Taxas específicas de mortalidade por sexo e grupo etário. Amapá, anos selecionados.}
    \includegraphics[width=0.95\textwidth]{imagens/nmx_painel_ap.png}
    \label{fig:nmx_painel_q3}
\end{figure}

\subsection{Mortalidade Infantil e Perinatal}

A Taxa de Mortalidade Infantil (TMI) foi calculada conforme solicitado no enunciado: média dos óbitos infantis de 2022 a 2024 no numerador e nascidos vivos de 2023 no denominador.

\[
TMI = \frac{\overline{D}_{<1,\,2022-2024}}{NV_{2023}} \times 1000
\]

\begin{table}[H]
\centering
\small
\caption{Taxa de Mortalidade Infantil, total e por sexo. Amapá.}
\label{tab:tmi_q3}
\begin{tabular}{lrrr}
\toprule
\textbf{Sexo} & \textbf{Média de óbitos infantis} & \textbf{NV 2023} & \textbf{TMI por mil} \\
\midrule
Total & 238,3 & 12.948 & 18,41 \\
Masculino & 128,7 & 6.637 & 19,39 \\
Feminino & 107,7 & 6.310 & 17,06 \\
\bottomrule
\end{tabular}
\end{table}

A TMI masculina foi maior que a feminina, padrão comum em estudos de mortalidade infantil. A decomposição mostra maior peso dos óbitos neonatais precoces, isto é, ocorridos de 0 a 6 dias de vida.

\begin{table}[H]
\centering
\small
\caption{Componentes da mortalidade infantil. Amapá.}
\label{tab:componentes_tmi_q3}
\resizebox{\textwidth}{!}{%
\begin{tabular}{lrr}
\toprule
\textbf{Componente} & \textbf{Média de óbitos 2022--2024} & \textbf{Taxa por mil} \\
\midrule
Neonatal precoce (0--6 dias) & 112,3 & 8,68 \\
Neonatal tardio (7--27 dias) & 32,0 & 2,47 \\
Neonatal total (0--27 dias) & 144,3 & 11,15 \\
Pós-neonatal (28 dias--menos de 1 ano) & 94,0 & 7,26 \\
\bottomrule
\end{tabular}
}
\end{table}

A mortalidade perinatal foi calculada somando os óbitos fetais de 28 semanas ou mais aos óbitos neonatais precoces. Também foi calculada uma medida de sensibilidade com todos os óbitos fetais, pois a variável de idade gestacional pode ter preenchimento incompleto.

\[
TMP = \frac{OF_{28+} + D_{0-6}}{NV + OF_{28+}} \times 1000
\]

em que $OF_{28+}$ representa os óbitos fetais com 28 semanas ou mais de gestação e $D_{0-6}$ representa os óbitos neonatais precoces.

\begin{table}[H]
\centering
\small
\caption{Taxa de Mortalidade Perinatal. Amapá.}
\label{tab:perinatal_q3}
\resizebox{\textwidth}{!}{%
\begin{tabular}{lrrrr}
\toprule
\textbf{Critério} & \textbf{Fetais} & \textbf{Neo. precoces} & \textbf{Denominador} & \textbf{TMP por mil} \\
\midrule
Fetais 28+ semanas & 97,7 & 112,3 & 13.059 & 16,08 \\
Todos os fetais -- sensibilidade & 148,7 & 112,3 & 13.108 & 19,91 \\
\bottomrule
\end{tabular}
}
\end{table}

Os resultados indicam que a mortalidade infantil ainda tem peso importante no componente neonatal, principalmente nos primeiros dias de vida. Isso sugere relação com condições da gestação, do parto, do nascimento e da assistência neonatal.

\subsection{Estrutura de Mortalidade por Causas}

A análise por causas utilizou a causa básica do óbito, agrupada segundo capítulos da CID-10. A Tabela \ref{tab:causas_ano_q3} apresenta os principais grupos de causas em 2010, 2021 e 2024.

\begin{table}[H]
\centering
\scriptsize
\caption{Principais grupos de causas de morte. Amapá, 2010, 2021 e 2024.}
\label{tab:causas_ano_q3}
\resizebox{\textwidth}{!}{%
\begin{tabular}{llrr}
\toprule
\textbf{Ano} & \textbf{Grupo de causa} & \textbf{Óbitos} & \textbf{\% do ano} \\
\midrule
2010 & Causas externas & 522 & 22,9 \\
2010 & Aparelho circulatório & 367 & 16,1 \\
2010 & Sintomas e achados anormais & 297 & 13,0 \\
2010 & Neoplasias & 235 & 10,3 \\
2010 & Aparelho respiratório & 214 & 9,4 \\
\midrule
2021 & COVID-19 (B34.2/U07) & 1.033 & 22,6 \\
2021 & Causas externas & 756 & 16,5 \\
2021 & Aparelho circulatório & 700 & 15,3 \\
2021 & Neoplasias & 501 & 10,9 \\
2021 & Aparelho respiratório & 343 & 7,5 \\
\midrule
2024 & Aparelho circulatório & 924 & 24,3 \\
2024 & Causas externas & 684 & 18,0 \\
2024 & Neoplasias & 572 & 15,1 \\
2024 & Aparelho respiratório & 451 & 11,9 \\
2024 & Endócrinas, nutricionais e metabólicas & 214 & 5,6 \\
\bottomrule
\end{tabular}
}
\end{table}

Em 2010, as causas externas foram o principal grupo de causas, com forte peso entre homens e idades jovens. Em 2021, a COVID-19 aparece como principal causa agregada, representando 22,6\% dos óbitos. Em 2024, o perfil volta a ser mais marcado por doenças crônicas, com destaque para doenças do aparelho circulatório, neoplasias e doenças respiratórias.

\begin{figure}[H]
    \centering
    \caption{Estrutura de mortalidade por grupos de causas. Amapá, 2010, 2021 e 2024.}
    \includegraphics[width=0.95\textwidth]{imagens/causas_top20_ano_ap.png}
    \label{fig:causas_ano_q3}
\end{figure}

A análise por sexo mostra diferenças importantes. Entre os homens, as causas externas têm participação elevada, especialmente em 2010 e 2024. Entre as mulheres, há maior peso relativo das doenças do aparelho circulatório, neoplasias e doenças respiratórias.

\begin{figure}[H]
    \centering
    \caption{Estrutura de causas de morte por sexo. Amapá, anos selecionados.}
    \includegraphics[width=0.95\textwidth]{imagens/causas_top20_sexo_ap.png}
    \label{fig:causas_sexo_q3}
\end{figure}

A análise por grupo de idade também mostrou diferenças importantes na estrutura de causas. Em menores de 5 anos, predominaram causas relacionadas ao período perinatal. Nos grupos de 5 a 14 anos e 15 a 39 anos, as causas externas tiveram maior peso relativo, especialmente entre jovens e adultos. A partir de 40 anos, aumentou a participação das doenças crônicas, principalmente doenças do aparelho circulatório, neoplasias e doenças respiratórias. Em 2021, a COVID-19 teve participação importante nos grupos de 40 a 59 anos e 60 anos ou mais.

\begin{table}[H]
\centering
\scriptsize
\caption{Principal grupo de causa de morte por idade. Amapá, 2010, 2021 e 2024.}
\label{tab:causas_idade_q3}
\resizebox{\textwidth}{!}{%
\begin{tabular}{lllrr}
\toprule
\textbf{Ano} & \textbf{Grupo de idade} & \textbf{Principal causa} & \textbf{Óbitos} & \textbf{\%} \\
\midrule
2010 & $<5$   & Período perinatal & 213 & 61,9 \\
2010 & 5--14  & Causas externas & 29 & 44,6 \\
2010 & 15--39 & Causas externas & 350 & 63,5 \\
2010 & 40--59 & Causas externas & 98 & 22,6 \\
2010 & 60+    & Aparelho circulatório & 246 & 27,9 \\
2021 & $<5$   & Período perinatal & 158 & 48,2 \\
2021 & 5--14  & Causas externas & 27 & 42,9 \\
2021 & 15--39 & Causas externas & 521 & 59,4 \\
2021 & 40--59 & COVID-19 & 340 & 31,8 \\
2021 & 60+    & COVID-19 & 603 & 26,9 \\
2024 & $<5$   & Período perinatal & 107 & 41,0 \\
2024 & 5--14  & Causas externas & 11 & 24,4 \\
2024 & 15--39 & Causas externas & 442 & 65,0 \\
2024 & 40--59 & Aparelho circulatório & 193 & 24,9 \\
2024 & 60+    & Aparelho circulatório & 688 & 33,8 \\
\bottomrule
\end{tabular}
}
\end{table}

A mortalidade por COVID-19 concentrou-se em 2021, principalmente entre pessoas de 60 anos ou mais. Também houve maior número absoluto de óbitos entre homens. Em 2024, os óbitos por COVID-19 foram residuais. O registro isolado em 2010 não deve ser interpretado como COVID-19 pandêmica, mas como registro associado aos códigos usados no agrupamento.

\begin{figure}[H]
    \centering
    \caption{Óbitos por COVID-19 segundo sexo e grupo de idade. Amapá, 2010, 2021 e 2024.}
    \includegraphics[width=0.9\textwidth]{imagens/covid_sexo_idade_ap.png}
    \label{fig:covid_q3}
\end{figure}

\subsection{Tábuas de Vida}

Foram construídas tábuas de vida abreviadas para homens e mulheres em 2010 e 2024. A tábua de vida parte das taxas específicas de mortalidade e estima funções como $l_x$, ${}_nq_x$, ${}_nL_x$, $T_x$ e $e_x^0$.

Para as primeiras idades, foram usados fatores de separação estimados com os microdados do SIM. Para as demais idades, adotou-se a hipótese usual de distribuição aproximadamente uniforme dos óbitos no intervalo.

Na construção da tábua de vida, as taxas foram utilizadas em forma decimal, isto é, sem multiplicação por mil. Assim, foi considerada:

\[
{}_nm_x = \frac{{}_nD_x}{{}_nP_x}
\]

A relação utilizada entre ${}_nm_x$ e ${}_nq_x$ foi:

\[
{}_nq_x =
\frac{n \cdot {}_nm_x}{1 + (n-a_x)\cdot {}_nm_x}
\]

em que ${}_nm_x$ representa a taxa específica de mortalidade em forma decimal e $a_x$ representa o tempo médio vivido no intervalo pelos indivíduos que morreram naquele intervalo.

\begin{table}[H]
\centering
\small
\caption{Fatores de separação estimados para as primeiras idades. Amapá.}
\label{tab:fatores_sep_q3}
\resizebox{\textwidth}{!}{%
\begin{tabular}{llrrrr}
\toprule
\textbf{Ano} & \textbf{Sexo} & \textbf{$a_0$} & \textbf{$a_{1-4}$} & \textbf{Óbitos 0} & \textbf{Óbitos 1--4} \\
\midrule
2010 & Feminino & 0,100 & 1,595 & 144 & 21 \\
2010 & Masculino & 0,090 & 1,833 & 155 & 24 \\
2024 & Feminino & 0,086 & 1,350 & 107 & 20 \\
2024 & Masculino & 0,114 & 1,667 & 108 & 24 \\
\bottomrule
\end{tabular}
}
\end{table}

A Tabela \ref{tab:resumo_tv_q3} apresenta a esperança de vida ao nascer e aos 60 anos. Em todos os casos, a esperança de vida feminina foi maior que a masculina.

\begin{table}[H]
\centering
\small
\caption{Esperança de vida ao nascer e aos 60 anos pelas tábuas calculadas. Amapá.}
\label{tab:resumo_tv_q3}
\begin{tabular}{llrr}
\toprule
\textbf{Ano} & \textbf{Sexo} & \textbf{$e_0^0$} & \textbf{$e_{60}^0$} \\
\midrule
2010 & Masculino & 75,20 & 24,95 \\
2010 & Feminino & 83,14 & 28,35 \\
2024 & Masculino & 71,83 & 20,75 \\
2024 & Feminino & 79,72 & 24,89 \\
\bottomrule
\end{tabular}
\end{table}

A comparação com a tábua do IBGE mostra que os valores calculados para 2024 ficaram próximos aos valores da Revisão 2024. Para 2010, as estimativas calculadas ficaram acima das estimativas do IBGE, o que exige cautela. Essa diferença pode estar relacionada à qualidade do registro de óbitos, à estabilidade das taxas em populações menores e ao uso direto das taxas observadas do SIM.

\begin{table}[H]
\centering
\scriptsize
\caption{Comparação da esperança de vida calculada com a tábua do IBGE. Amapá.}
\label{tab:comp_ibge_q3}
\resizebox{\textwidth}{!}{%
\begin{tabular}{llrrrr}
\toprule
\textbf{Ano} & \textbf{Sexo} & \textbf{$e_0$ calc.} & \textbf{$e_0$ IBGE} & \textbf{$e_{60}$ calc.} & \textbf{$e_{60}$ IBGE} \\
\midrule
2010 & Masculino & 75,20 & 70,81 & 24,95 & 20,39 \\
2010 & Feminino & 83,14 & 78,58 & 28,35 & 24,24 \\
2024 & Masculino & 71,83 & 70,26 & 20,75 & 20,44 \\
2024 & Feminino & 79,72 & 78,78 & 24,89 & 23,45 \\
\bottomrule
\end{tabular}
}
\end{table}

\begin{figure}[H]
    \centering
    \caption{Função $l_x$ da tábua de vida. Amapá, 2010 e 2024.}
    \includegraphics[width=0.78\textwidth]{imagens/tv_lx_ap.png}
    \label{fig:lx_q3}
\end{figure}

\begin{figure}[H]
    \centering
    \caption{Função ${}_nq_x$ da tábua de vida. Amapá, 2010 e 2024.}
    \includegraphics[width=0.78\textwidth]{imagens/tv_nqx_ap.png}
    \label{fig:nqx_q3}
\end{figure}

A função $l_x$ mostra maior sobrevivência feminina ao longo das idades. A função ${}_nq_x$ apresenta o padrão esperado: maior probabilidade de morte no primeiro ano de vida, queda nas idades infantis e aumento progressivo nas idades adultas e idosas. As probabilidades masculinas são superiores às femininas em grande parte da vida adulta.

\subsection{Comentários}

Os resultados da Questão 3 são coerentes com o padrão esperado de mortalidade. A mortalidade aumenta com a idade, é maior entre homens em grande parte dos grupos etários e apresenta forte peso das causas externas entre homens jovens e adultos.

O ano de 2021 se destaca pelo impacto da COVID-19, tanto no aumento da TBM quanto na estrutura de causas de morte. Em 2024, observa-se retorno a um perfil mais associado às doenças crônicas, principalmente doenças do aparelho circulatório, neoplasias e doenças respiratórias.

As tábuas de vida reforçam o diferencial por sexo, com maior esperança de vida feminina. Os resultados de 2024 devem ser interpretados com cautela, pois os dados do DATASUS podem ser preliminares. Além disso, diferenças entre a tábua calculada e a tábua do IBGE mostram a importância da qualidade do SIM e da estabilidade das taxas em populações menores, como é o caso do Amapá.

